\begin{frame}[t]{Combining branch results should add information while accounting for overlap.}
\lead{\textbf{Fan-out} explores alternatives. \textbf{Fan-in} turns their results into a useful decision.}
\vspace{.25cm}
\lead{\textbf{Choose or merge artifacts:} check the final code, model or candidate.}
\vspace{.15cm}
\lead{\textbf{Pool measurements:} use their reliability and shared information.}
\vspace{.15cm}
\lead{\textbf{Goal:} reduce uncertainty without counting the same information twice.}
\sources{Next: merge rules, overlapping observations, statistical weights and the resulting precision.}

\qadetail{This bridge states the responsibilities without repeating the component-architecture diagram. Two private executions can use the same test fixture, prompt, grader, remote response or recorded observation. Artifact selection, exact composition and common-target measurement pooling have different checks, as the following menu explains. The integrated external recorder and abstaining evaluator are proposed; the built reducer carries declared evidence identity and lineage under its stated statistical model.}
\note{[0:40] We have seen how to create and control alternative runs. Now comes the central fan-in question: how do we turn their results into a better decision? For code or model candidates, we select or compose an artifact and check it. For measurements of one quantity, we combine information using reliability and overlap. The aim is a clearer signal and less uncertainty, with shared information accounted for. We will distinguish these merge rules, then show an overlap diagram and the statistical calculation.}
\end{frame}
